% Generated draft: requires mathematical and source review.
\documentclass[11pt]{article}
\usepackage[margin=25mm]{geometry}
\usepackage{fontspec}
\setmainfont{Latin Modern Roman}
\usepackage{amsmath,amssymb,mathtools,unicode-math}
\setmathfont{Latin Modern Math}
\usepackage{enumitem,needspace,xurl,hyperref}
\setlength{\parindent}{0pt}
\setlength{\parskip}{5pt}
\setlength{\emergencystretch}{3em}
\newcommand{\N}{\mathbb{N}}
\newcommand{\Z}{\mathbb{Z}}
\newcommand{\Q}{\mathbb{Q}}
\newcommand{\R}{\mathbb{R}}
\newcommand{\C}{\mathbb{C}}
\newcommand{\F}{\mathbb{F}}
\newcommand{\A}{\mathbb{A}}
\newcommand{\PP}{\mathbb{P}}
\newcommand{\E}{\mathbb{E}}
\newcommand{\eps}{\varepsilon}
\newcommand{\dd}{\,\mathrm{d}}
\begin{document}
\section*{UnsolvedMath ID: 20001960}
\textbf{Source problem code:} catalogue entry
\subsection{Definitions and mathematical statement}
E. Lerman's Question \(5.8\) asks for a proof of Kirwan surjectivity without Morse theory. In the usual compact Hamiltonian setting, with compact acting group \(G\), moment map \(\mu\), and an applicable regular level, the target map is restriction
\[
 H_G^*(X;\Q)\longrightarrow H_G^*(\mu^{-1}(0);\Q),
\]
with quotient interpretation when its hypotheses permit. The source does not restate a complete version of the symplectic theorem. It motivates an alternative method by [CDM88]'s non-Morse connectedness and convexity proofs and their usefulness for contact moment maps or symplectic cones. It notes the failure of the usual critical-point/isotropy and properness ingredients there. This motivation does not assert that ordinary contact Kirwan surjectivity is always true.

The aid's non-Morse extension criterion uses only equivariant Mayer--Vietoris. Let \(Z\subset X\) be invariant, choose invariant open sets \(U,V\) with \(U\cup V=X\), and assume \(U\) equivariantly deformation retracts to \(Z\). Write \(W=U\cap V\). For \(c\in H_G^k(Z;R)\), let \(a\in H_G^k(U;R)\) be its corresponding class. Define
\[
 \operatorname{obs}(c)=[a|_W]\ \in\
 H_G^k(W;R)/\operatorname{im}(H_G^k(V;R)\to H_G^k(W;R)).
\]
Exactness gives
\[
 c\in\operatorname{im}(H_G^k(X;R)\to H_G^k(Z;R))
 \quad\Longleftrightarrow\quad \operatorname{obs}(c)=0.
\]
Thus this cokernel detects precisely the extension obstruction for the chosen cover. It is not a proof that the obstruction always vanishes for a moment map.

For a contact manifold \((M,\alpha)\) with a strict contact \(G\)-action, radial symplectization is
\[
 S=M\times\R_{>0},\qquad \omega=\dd(r\alpha),\qquad
 \mu_S(m,r)=r\mu_M(m).
\]
At zero, \(\mu_S^{-1}(0)=\mu_M^{-1}(0)\times\R_{>0}\), and radial projection gives equivariant homotopy equivalences of the full space and its zero level. Hence the contact and symplectization zero-level restriction maps are conjugate by cohomology isomorphisms. This assertion is specific to zero; it supplies no properness or surjectivity by itself.

The aid's explicit obstruction is Lerman's standard contact \(S^3\) with action
\[
 \lambda(z_1,z_2)=(\lambda z_1,\lambda^{-1}z_2),\qquad
 \mu_M=|z_1|^2-|z_2|^2.
\]
Its image \([-1,1]\) is convex, every fiber is connected, and zero is regular. For a coefficient field \(k\), the Borel map is
\[
 k[u]/(u^2)\longrightarrow\bigwedge_k(\eta),\qquad
 |u|=2,\ |\eta|=1,\quad u\longmapsto0,
\]
where \(\eta^2=0\). Its kernel is \((u)\) and its graded-module cokernel is \(k\eta\) in degree one. The nonextendable class survives radial symplectization. Therefore convexity and connectedness alone do not imply ordinary contact surjectivity, and a successful alternative proof must use additional hypotheses particular to its intended symplectic setting. The source's following Comment \(5.9\) is cut off after ``As mentioned in''; no missing claim is supplied.
\subsection{Short English statement}
The source asks for a non-Morse proof of Kirwan surjectivity. The aid gives an exact Mayer--Vietoris test for extending zero-level classes and shows the zero-level map is preserved by radial symplectization. Lerman's sphere still has a nonextendable degree-one class despite convex moment image and connected fibers, so those properties alone cannot prove a contact analogue.
\subsection{Sources}
\begin{itemize}
\item Local UnsolvedMath registry, permanent ID \(20001960\), source code \texttt{catalogue entry}; canonical collection \texttt{aim-geometry-notes.json}, notes[297]. Statement status \texttt{reconstructed\_unverified}; clean statement null.
\item AIM workshop ``Moment maps and surjectivity in various geometries,'' Question \(5.8\), E. Lerman, and truncated Comment \(5.9\): \url{https://aimath.org/WWN/momentmaps/momentmaps.pdf}. No PDF was opened in this drafting session.
\item [CDM88] is the supplied reference for the alternative connectedness/convexity approach; no full bibliographic entry is given. Lerman's contact-sphere example is used by the aid. Original remarks and literature fields are empty, with no proposer-field value or proposed year.
\item Supplied research attempt \(1\), \texttt{partial\_result}, run status \texttt{partially\_solved}, rigor label ``proved,'' novelty confidence ``low.'' Registry classification \texttt{PARTIAL-PROGRESS}, updated 14 August 2026.
\end{itemize}
\end{document}
